\documentclass{article}
\usepackage[a4paper, margin=2cm]{geometry}
\usepackage{bussproofs}
\usepackage{extarrows}
\usepackage{mathtools}
\usepackage{amssymb}
\usepackage{parskip}
\usepackage{multicol}
\usepackage{amsthm}

\newcommand{\enc}{[\![}
\newcommand{\cne}{]\!]}
\newcommand{\N}{\mathbb{N}}
\renewcommand{\arraystretch}{2}


\begin{document}

% ================================================================
% CCS SEMANTICS
% ================================================================

% ---------------------------------------------------------------
\section*{Encoding $\enc\cdot\cne : \text{CCS-VP} \to \text{CCS}$}
% ---------------------------------------------------------------

\[
    \begin{array}{rcll}
        \enc a(x).P \cne                       & = & \sum_{n \in \N} a_n .  \enc  P\{^n/_x\} \cne                  & \\
        \enc \bar{a}(e).P \cne                 & = & \bar{a}_m .  \enc  P \cne \quad                               & e \downarrow m \\
        \enc \tau.P \cne                       & = & \tau .  \enc P \cne                                           & \\
        \enc P \mid Q \cne                     & = & \enc P \cne  \mid  \enc Q \cne                                & \\
        \enc \sum_{i \in I} P_i \cne           & = & \displaystyle\sum_{i \in I}  \enc P_i \cne                    & \\
        \enc P \setminus L \cne                & = & \enc P \cne  \setminus \{a_m \mid a \in L, m \in \mathbb{N}\} & \\[6pt]
        \enc P[f \cne ]                        & = & \enc P \cne [f'] \quad f'(a_m) = f(a)_m \\[6pt]
        \enc \text{if } b \text{ then } P \cne & = & 
            \begin{cases} 
                \enc P \cne & b \text{ true} \\
                0           & \text{otherwise} 
            \end{cases} \\
        \enc K(e_1,\ldots,e_h) \cne            & = & K_{n_1,\ldots,n_h}                                            & e_i \downarrow n_i \\
        \\
        K_{n_1,\ldots,n_h}                     & \overset{\text{def}}{=} & \enc P\{^{n_1}/_{x_1},\ldots,^{n_h}/_{x_h}\} \cne & \\
                                               & & \text{where } K(x_1,\ldots,x_h) \overset{\text{def}}{=} P         &\\ 
    \end{array}
\]

\newpage

% ================================================================
% ENCODING
% ================================================================

% ---------------------------------------------------------------
\section*{CCS Operational Semantics}
% ---------------------------------------------------------------
\vspace{1cm}
  % ACT
  \begin{prooftree}
    \AxiomC{}
    \RightLabel{\textsc{Act}}
    \UnaryInfC{$\alpha.P \xrightarrow{\alpha} P$}
  \end{prooftree}

  % SUM
  \begin{prooftree}
    \AxiomC{$P_j \xrightarrow{\alpha} P_j'$}
    \RightLabel{\textsc{Sum} \quad $j \in I$}
    \UnaryInfC{$\displaystyle\sum_{i \in I} P_i \xrightarrow{\alpha} P_j'$}
  \end{prooftree}

  % PAR-L
  \begin{prooftree}
    \AxiomC{$P \xrightarrow{\alpha} P'$}
    \RightLabel{\textsc{Par-L}}
    \UnaryInfC{$P \mid Q \xrightarrow{\alpha} P' \mid Q$}
  \end{prooftree}

  % PAR-R
  \begin{prooftree}
    \AxiomC{$Q \xrightarrow{\alpha} Q'$}
    \RightLabel{\textsc{Par-R}}
    \UnaryInfC{$P \mid Q \xrightarrow{\alpha} P \mid Q'$}
  \end{prooftree}

  % COM
    \begin{prooftree}
      \AxiomC{$P \xrightarrow{\alpha} P'$}
      \AxiomC{$Q \xrightarrow{\bar{\alpha}} Q'$}
      \RightLabel{\textsc{Com}}
      \BinaryInfC{$P \mid Q \xrightarrow{\tau} P' \mid Q'$}
    \end{prooftree}

  % RES
  \begin{prooftree}
    \AxiomC{$P \xrightarrow{\alpha} P'$}
    \RightLabel{\textsc{Res} \quad $\alpha, \bar{\alpha} \notin L$}
    \UnaryInfC{$P \setminus L \xrightarrow{\alpha} P' \setminus L$}
  \end{prooftree}
  
  % RENAME
  \begin{prooftree}
    \AxiomC{$P \xrightarrow{\alpha} P'$}
    \RightLabel{\textsc{Ren}}
    \UnaryInfC{$P[f] \xrightarrow{f(\alpha)} P'[f]$}
  \end{prooftree}

  % CONST
    \begin{prooftree}
      \AxiomC{$P \xrightarrow{\alpha} P'$}
      \RightLabel{\textsc{Const} \quad $K \overset{\text{def}}{=} P$}
      \UnaryInfC{$K \xrightarrow{\alpha} P'$}
    \end{prooftree}

% ================================================================
% THEOREM
% ================================================================

For all $P$ in CCS-VP
\\[2em]
\[\text{(i)}\quad P \xrightarrow{\alpha} P' \;\Longrightarrow\; \enc P\cne \xrightarrow{\hat\alpha} \enc P'\cne\]
\[\text{(ii)}\quad \enc P\cne \xrightarrow{\hat\alpha} Q \;\Longrightarrow\; \exists P'.\; P \xrightarrow{\alpha} P' \;\wedge\; \enc P'\cne = Q\]
\\[2em]
\(\text{where }\;\hat\alpha = \begin{cases} a_m & \alpha = a(m) \\ \bar{a}_m & \alpha = \bar{a}(m) \\ \tau & \alpha = \tau \end{cases}\)
\\[2em]

% ================================================================
% PROOF OVERVIEW
% ================================================================

Both implications are established by \emph{rule induction}, or induction on the depth of
the derivation tree, which is the natural choice here because the transition relation
is itself defined inductively by inference rules. The natural alternative would be 
\emph{structural induction on $P$}, case-splitting on its syntactic form rather than 
on the last rule fired. For most forms this coincides with rule induction, but it fails 
cleanly at \textsc{Const}: a constant $K$ is a syntactic atom with no sub-term, so structural
induction provides no smaller process to apply the IH to. Rule induction avoids this by
measuring derivation depth, which strictly decreases at every premise, including the 
unfolding step inside \textsc{Const}.

For (i) we proceed by induction on the derivation of $P \xrightarrow{\alpha} P'$, with a case
analysis on the last rule applied. The \emph{base cases} are the three axiom forms 
of CCS-VP (input, output, and $\tau$ prefix): for each, the encoding $\enc P \cne$ is 
a concrete CCS term that admits the required transition directly via the Act or Sum rule.
For the \emph{inductive cases} (Sum, Par-L, Par-R, Com, Res, Ren, Const, Cond), the
induction hypothesis supplies transitions for the shorter sub-derivation(s), and the
matching CCS rule then closes the goal.

For (ii) we instead proceed by induction on the derivation of $\enc P \cne \xrightarrow{\hat\alpha} Q$,
with a case analysis on the syntactic form of $P$. For each form of $P$, the structure of 
$\enc P \cne$ uniquely determines which CCS rule could have fired last; the induction 
hypothesis then recovers a CCS-VP witness $P'$ with $P \xrightarrow{\alpha} P'$, and the 
encoding equations confirm $\enc P' \cne = Q$.

Throughout both directions the key invariant is that $\hat\alpha$ correctly bridges
CCS-VP labels and CCS channels $a(m)\leftrightarrow a_m$,\; $\bar{a}(m)\leftrightarrow\bar{a}_m$,\; $\tau\leftrightarrow\tau$
ensuring that complementary actions in CCS-VP correspond to complementary channels
in CCS, so that synchronization is preserved by the encoding.

\newpage

% ================================================================
% PROOF (i)
% ================================================================

\section*{\centering \(P \xrightarrow{\alpha} P' \;\Longrightarrow\; \enc P\cne \xrightarrow{\hat\alpha} \enc P'\cne\)}
\vspace{.5cm}
\subsubsection*{\centering\emph{Base cases}}
\vspace{.5cm}
\textbf{Input} $a(x).P \xrightarrow{a(m)} P\{^m/_x\}$:\\[1em]
\begin{prooftree}
    \AxiomC{$a_m.\enc P\{^m/_x\}\cne \xrightarrow{a_m} \enc P\{^m/_x\}\cne$}
    \RightLabel{Sum}
    \UnaryInfC{$\displaystyle\sum_{n \in \N} a_n.\enc P\{^n/_x\}\cne \xrightarrow{a_m} \enc P\{^m/_x\}\cne$}
    \RightLabel{=}
    \UnaryInfC{$\enc a(x).P\cne \xrightarrow{a_m} \enc P\{^m/_x\}\cne $}
\end{prooftree}

\textbf{Output} $\bar{a}(e).P \xrightarrow{\bar{a}(m)} P$, $e\downarrow m$:\\[1em]
\begin{prooftree}
    \AxiomC{}
    \RightLabel{Act}
    \UnaryInfC{$ \bar{a}_m.\enc P \cne \xrightarrow{\bar{a}_m} \enc P\cne$}
    \RightLabel{=}
    \UnaryInfC{$\enc \bar{a}(e).P \cne \xrightarrow{\bar{a}_m} \enc P \cne $}
\end{prooftree}

\textbf{Silent} $\tau.P \xrightarrow{\tau} P$:\\[1em]
\begin{prooftree}
    \AxiomC{}
    \RightLabel{Act}
    \UnaryInfC{$ \tau.\enc P \cne \xrightarrow{\tau} \enc P\cne$}
    \RightLabel{=}
    \UnaryInfC{$\enc \tau.P \cne \xrightarrow{\tau} \enc P \cne $}
\end{prooftree}

\vspace{.5cm}
\subsubsection*{\centering\emph{Inductive cases}}

\begin{prooftree}
    \AxiomC{$P_i \xrightarrow{\alpha} P_i'$}
    \RightLabel{IH}
    \UnaryInfC{$\enc P_i \cne \xrightarrow{\hat\alpha} \enc P_i' \cne $}
\end{prooftree}

\vspace{.5cm}

\textbf{Sum} $\displaystyle\sum_{i \in I} P_i \xrightarrow{\alpha} P_j' \qquad j \in I$:\\[1em]
\begin{prooftree}
    \AxiomC{$P_j \xrightarrow{\alpha} P_j'$}
    \RightLabel{IH}
    \UnaryInfC{$\enc P_j\cne \xrightarrow{\hat\alpha} \enc P_j'\cne$}
    \RightLabel{Sum}
    \UnaryInfC{$\displaystyle\sum_{i \in I}\enc P_i\cne \xrightarrow{\hat\alpha} \enc P_j'\cne$}
    \RightLabel{$=$}
    \UnaryInfC{$\enc\displaystyle\sum_{i \in I} P_i\cne \xrightarrow{\hat\alpha} \enc P_j'\cne$}
\end{prooftree}

\textbf{Par-L}: $P_1\mid P_2 \xrightarrow{\alpha} P_1'\mid P_2$ \\[1em]
\begin{prooftree}
  \AxiomC{$P_1 \xrightarrow{\alpha} P_1'$}
  \RightLabel{IH}
  \UnaryInfC{$\enc P_1\cne \xrightarrow{\hat\alpha} \enc P_1'\cne$}
  \RightLabel{Par-L}
  \UnaryInfC{$\enc P_1\cne \mid \enc P_2\cne \xrightarrow{\hat\alpha} \enc P_1'\cne \mid \enc P_2\cne$}
  \RightLabel{$=$}
  \UnaryInfC{$\enc P_1 \mid P_2\cne \xrightarrow{\hat\alpha} \enc P_1' \mid P_2\cne$}
\end{prooftree}

\textbf{Par-R}: $P_1\mid P_2 \xrightarrow{\alpha} P_1\mid P_2'$ \\[1em]
\begin{prooftree}
  \AxiomC{$P_2 \xrightarrow{\alpha} P_2'$}
  \RightLabel{IH}
  \UnaryInfC{$\enc P_2\cne \xrightarrow{\hat\alpha} \enc P_2'\cne$}
  \RightLabel{Par-R}
  \UnaryInfC{$\enc P_1\cne \mid \enc P_2\cne \xrightarrow{\hat\alpha} \enc P_1\cne \mid \enc P_2'\cne$}
  \RightLabel{$=$}
  \UnaryInfC{$\enc P_1 \mid P_2\cne \xrightarrow{\hat\alpha} \enc P_1 \mid P_2'\cne$}
\end{prooftree}

\textbf{Com} $P_1 \mid P_2 \xrightarrow{\tau} P_1' \mid P_2'$:\\[1em]
\begin{prooftree}
    \AxiomC{$P_1 \xrightarrow{a(m)} P_1'$}
    \RightLabel{IH}
    \UnaryInfC{$\enc P_1\cne \xrightarrow{a_m} \enc P_1'\cne$}
    \AxiomC{$P_2 \xrightarrow{\bar{a}(m)} P_2'$}
    \RightLabel{IH}
    \UnaryInfC{$\enc P_2\cne \xrightarrow{\bar{a}_m} \enc P_2'\cne$}
    \RightLabel{Com}
    \BinaryInfC{$\enc P_1\cne \mid \enc P_2\cne \xrightarrow{\tau} \enc P_1'\cne \mid \enc P_2'\cne$}
    \RightLabel{$=$}
    \UnaryInfC{$\enc P_1 \mid P_2\cne \xrightarrow{\tau} \enc P_1' \mid P_2'\cne$}
\end{prooftree}

\textbf{Res} $P\setminus L \xrightarrow{\alpha} P'\setminus L$, $\alpha\notin L\cup\bar{L}$:\\[1em]
\begin{prooftree}
    \AxiomC{$P \xrightarrow{\alpha} P'$}
    \RightLabel{IH}
    \UnaryInfC{$\enc P\cne \xrightarrow{\hat\alpha} \enc P'\cne$}
    \RightLabel{Res}
    \UnaryInfC{$\enc P\cne \setminus \{a_m \mid a \in L,\,m \in \N\} \xrightarrow{\hat\alpha} \enc P'\cne \setminus \{a_m \mid a \in L,\,m \in \N\}$}
    \RightLabel{$=$}
    \UnaryInfC{$\enc P\setminus L\cne \xrightarrow{\hat\alpha} \enc P'\setminus L\cne$}
\end{prooftree}

\textbf{Ren} $P[f] \xrightarrow{f(\alpha)} P'[f]$, where $f'(a_m)=f(a)_m$ and $\widehat{f(\alpha)}=f'(\hat\alpha)$:\\[1em]
\begin{prooftree}
    \AxiomC{$P \xrightarrow{\alpha} P'$}
    \RightLabel{IH}
    \UnaryInfC{$\enc P\cne \xrightarrow{\hat\alpha} \enc P'\cne$}
    \RightLabel{Ren}
    \UnaryInfC{$\enc P\cne[f'] \xrightarrow{f'(\hat\alpha)} \enc P'\cne[f']$}
    \RightLabel{$=$}
    \UnaryInfC{$\enc P[f]\cne \xrightarrow{\widehat{f(\alpha)}} \enc P'[f]\cne$}
\end{prooftree}

\textbf{Const} $K(e_1,\ldots,e_h) \xrightarrow{\alpha} R$, $e_i\downarrow n_i$:\\[1em]
\begin{prooftree}
    \AxiomC{$P\{^{n_1}/_{x_1},\ldots\} \xrightarrow{\alpha} R$}
    \RightLabel{IH}
    \UnaryInfC{$K_{n_1,\ldots} \xrightarrow{\hat\alpha} \enc R\cne$}
    \RightLabel{$=$}
    \UnaryInfC{$\enc K(e_1,\ldots,e_h)\cne \xrightarrow{\hat\alpha} \enc R\cne$}
\end{prooftree}

\textbf{Cond} $\text{if }b\text{ then }P \xrightarrow{\alpha} P'$, $b$ true:\\[1em]
\begin{prooftree}
    \AxiomC{$P \xrightarrow{\alpha} P'$}
    \RightLabel{IH}
    \UnaryInfC{$\enc P\cne \xrightarrow{\hat\alpha} \enc P'\cne$}
    \RightLabel{$=$}
    \UnaryInfC{$\enc\text{if }b\text{ then }P\cne \xrightarrow{\hat\alpha} \enc P'\cne$}
\end{prooftree}

\newpage

% ================================================================
% PROOF (ii)
% ================================================================

\section*{\centering\[\enc P\cne \xrightarrow{\hat\alpha} Q \;\Longrightarrow\; \exists P'.\; P \xrightarrow{\alpha} P' \;\wedge\; \enc P'\cne = Q\]}
\vspace{.5cm}
\subsubsection*{\centering\emph{Base cases}}
\vspace{.5cm}

\textbf{Input} [\;$a(x).P$\;]\\[1em]
\begin{prooftree}
    \AxiomC{}
    \RightLabel{Act}
    \UnaryInfC{$a(x).P \xrightarrow{a(m)} P\{^m/_x\}$}
\end{prooftree}
$\enc a(x).P\cne = \textstyle\sum_{n\in\N} a_n.\enc P\{^n/_x\}\cne$.\\
The only possible transition fires branch $a_m$ for some $m\in\N$, so $Q = \enc P\{^m/_x\}\cne$.\\
Taking $P' = P\{^m/_x\}$ witnesses $P \xrightarrow{a(m)} P'$ and $\enc P'\cne = Q$.
\vspace{.5cm}


\textbf{Output} [\;$\bar{a}(e).P$\;]\\[1em]
\begin{prooftree}
    \AxiomC{$e \downarrow m$}
    \RightLabel{Act}
    \UnaryInfC{$\bar{a}(e).P \xrightarrow{\bar{a}(m)} P$}
\end{prooftree}
$\enc\bar{a}(e).P\cne = \bar{a}_m.\enc P\cne$ since $e\downarrow m$.\\
The unique transition via Act forces $Q = \enc P\cne$.\\
Taking $P' = P$ witnesses $P \xrightarrow{\bar{a}(m)} P'$ and $\enc P'\cne = Q$.
\vspace{.5cm}

\textbf{Silent} [\;$\tau.P$\;]\\[1em]
\begin{prooftree}
    \AxiomC{}
    \RightLabel{Act}
    \UnaryInfC{$\tau.P \xrightarrow{\tau} P$}
\end{prooftree}
$\enc\tau.P\cne = \tau.\enc P\cne$.\\
The unique transition via Act forces $Q = \enc P\cne$.\\
Taking $P' = P$ witnesses $P \xrightarrow{\tau} P'$ and $\enc P'\cne = Q$.
\vspace{.5cm}

\newpage

\subsubsection*{\centering\emph{Inductive cases}}

\begin{prooftree}
    \AxiomC{$\enc P\cne \xrightarrow{\hat\alpha} Q$}
    \RightLabel{IH}
    \UnaryInfC{$P \xrightarrow{\alpha} P' \;\wedge\; \enc P'\cne = Q$}
\end{prooftree}

\vspace{.5cm}

\textbf{Sum} $\left[\;\displaystyle\sum_{i \in I} P_i\; \right]$\\[1em]
\begin{prooftree}
    \AxiomC{$\enc P_j\cne \xrightarrow{\hat\alpha} Q$}
    \RightLabel{IH}
    \UnaryInfC{$P_j \xrightarrow{\alpha} P_j' \;\wedge\; \enc P_j'\cne = Q$}
    \RightLabel{Sum}
    \UnaryInfC{$\displaystyle\sum_{i \in I} P_i \xrightarrow{\alpha} P_j'$}
\end{prooftree}
Since $\enc\textstyle\sum_{i\in I} P_i\cne = \textstyle\sum_{i\in I}\enc P_i\cne$, any transition must originate from some branch $j \in I$.\\
The IH recovers $P_j \xrightarrow{\alpha} P_j'$ with $\enc P_j'\cne = Q$.\\
Taking $P' = P_j'$ witnesses $P \xrightarrow{\alpha} P'$ and $\enc P'\cne = Q$.
\vspace{.5cm}

\textbf{Par} $\left[\;P_1\mid P_2\;\right]$\\[.5em]
$\enc P_1\mid P_2\cne = \enc P_1\cne\mid\enc P_2\cne$.\\ 
Any transition of $\enc P_1\cne\mid\enc P_2\cne$ must have arisen from one of three CCS rules.
\begin{itemize}
  \item \textsc{Par-L}
        \begin{prooftree}
            \AxiomC{$\enc P_1\cne \xrightarrow{\hat\alpha} Q_1$}
            \RightLabel{IH}
            \UnaryInfC{$P_1 \xrightarrow{\alpha} P_1' \;\wedge\; \enc P_1'\cne = Q_1$}
            \RightLabel{Par-L}
            \UnaryInfC{$P_1 \mid P_2 \xrightarrow{\alpha} P_1' \mid P_2$}
        \end{prooftree}
        The transition on the left component gives $Q = Q_1\mid\enc P_2\cne$.\\
        The IH recovers $P_1 \xrightarrow{\alpha} P_1'$ with $\enc P_1'\cne = Q_1$.\\
        Taking $P' = P_1'\mid P_2$ witnesses $P \xrightarrow{\alpha} P'$ and $\enc P'\cne = \enc P_1'\cne\mid\enc P_2\cne = Q_1\mid\enc P_2\cne = Q$.
  \item \textsc{Par-R}
        \begin{prooftree}
            \AxiomC{$\enc P_2\cne \xrightarrow{\hat\alpha} Q_2$}
            \RightLabel{IH}
            \UnaryInfC{$P_2 \xrightarrow{\alpha} P_2' \;\wedge\; \enc P_2'\cne = Q_2$}
            \RightLabel{Par-R}
            \UnaryInfC{$P_1 \mid P_2 \xrightarrow{\alpha} P_1 \mid P_2'$}
        \end{prooftree}
        The transition on the right component gives $Q = \enc P_1\cne\mid Q_2$.\\
        The IH recovers $P_2 \xrightarrow{\alpha} P_2'$ with $\enc P_2'\cne = Q_2$.\\
        Taking $P' = P_1\mid P_2'$ witnesses $P \xrightarrow{\alpha} P'$ and $\enc P'\cne = \enc P_1\cne\mid\enc P_2'\cne = \enc P_1\cne\mid Q_2 = Q$.
  \item \textsc{Com}
        \begin{prooftree}
            \AxiomC{$\enc P_1\cne \xrightarrow{a_m} Q_1$}
            \RightLabel{IH}
            \UnaryInfC{$P_1 \xrightarrow{a(m)} P_1' \;\wedge\; \enc P_1'\cne = Q_1$}
            \AxiomC{$\enc P_2\cne \xrightarrow{\bar{a}_m} Q_2$}
            \RightLabel{IH}
            \UnaryInfC{$P_2 \xrightarrow{\bar{a}(m)} P_2' \;\wedge\; \enc P_2'\cne = Q_2$}
            \RightLabel{Com}
            \BinaryInfC{$P_1 \mid P_2 \xrightarrow{\tau} P_1' \mid P_2'$}
        \end{prooftree}
        Complementary transitions $a_m$ and $\bar{a}_m$ map back to $a(m)$ and $\bar{a}(m)$.\\
        The IH on each side recovers $P_1 \xrightarrow{a(m)} P_1'$ with $\enc P_1'\cne = Q_1$ and $P_2 \xrightarrow{\bar{a}(m)} P_2'$ with $\enc P_2'\cne = Q_2$.\\
        Taking $P' = P_1'\mid P_2'$ witnesses $P \xrightarrow{\tau} P'$ and $\enc P'\cne = Q_1\mid Q_2 = Q$.
\end{itemize}
\vspace{.5cm}

\textbf{Res} $\left[\;P\setminus L\;\right]$\\[1em]
\begin{prooftree}
    \AxiomC{$\enc P\cne \xrightarrow{\hat\alpha} Q_0$}
    \RightLabel{IH}
    \UnaryInfC{$P \xrightarrow{\alpha} P' \;\wedge\; \enc P'\cne = Q_0$}
    \RightLabel{Res ($\alpha\notin L\cup\bar{L}$)}
    \UnaryInfC{$P \setminus L \xrightarrow{\alpha} P' \setminus L$}
\end{prooftree}
$\enc P\setminus L\cne = \enc P\cne\setminus\{a_m\mid a\in L,\,m\in\N\}$, so the transition decomposes as $\enc P\cne \xrightarrow{\hat\alpha} Q_0$ with $Q = Q_0\setminus\{a_m\mid a\in L,\,m\in\N\}$.\\
The transition being unblocked forces $\hat\alpha\notin\{a_m\mid a\in L\}$, hence $\alpha\notin L\cup\bar{L}$.\\
The IH recovers $P \xrightarrow{\alpha} P'$ with $\enc P'\cne = Q_0$.\\
Taking $P'\setminus L$ as the witness gives $P\setminus L \xrightarrow{\alpha} P'\setminus L$ and $\enc P'\setminus L\cne = \enc P'\cne\setminus\{a_m\mid a\in L\} = Q_0\setminus\{a_m\mid a\in L\} = Q$.
\vspace{.5cm}

\textbf{Ren} $\left[\;P[f]\;\right]$\\[1em]
\begin{prooftree}
    \AxiomC{$\enc P\cne \xrightarrow{\hat\alpha} Q_0$}
    \RightLabel{IH}
    \UnaryInfC{$P \xrightarrow{\alpha} P' \;\wedge\; \enc P'\cne = Q_0$}
    \RightLabel{Ren}
    \UnaryInfC{$P[f] \xrightarrow{f(\alpha)} P'[f]$}
\end{prooftree}
$\enc P[f]\cne = \enc P\cne[f']$, so the transition decomposes as $\enc P\cne \xrightarrow{\hat\alpha} Q_0$ with $Q = Q_0[f']$.\\
The IH recovers $P \xrightarrow{\alpha} P'$ with $\enc P'\cne = Q_0$.\\
Taking $P'[f]$ as the witness gives $P[f] \xrightarrow{f(\alpha)} P'[f]$ and $\enc P'[f]\cne = \enc P'\cne[f'] = Q_0[f'] = Q$.
\vspace{.5cm}

\textbf{Cond} $\left[\;\text{if }b\text{ then }P\;\right]$\\[1em]
\begin{prooftree}
    \AxiomC{$b \text{ true}$}
    \AxiomC{$\enc P\cne \xrightarrow{\hat\alpha} Q$}
    \RightLabel{IH}
    \UnaryInfC{$P \xrightarrow{\alpha} P' \;\wedge\; \enc P'\cne = Q$}
    \RightLabel{Cond}
    \BinaryInfC{$\text{if }b\text{ then }P \xrightarrow{\alpha} P'$}
\end{prooftree}
$\enc\text{if }b\text{ then }P\cne = \enc P\cne$ when $b$ is true.\\
The transition is directly on $\enc P\cne$.\\
The IH recovers $P \xrightarrow{\alpha} P'$ with $\enc P'\cne = Q$.\\
The IH witness $P'$ directly satisfies $P \xrightarrow{\alpha} P'$ and $\enc P'\cne = Q$.
\vspace{.5cm}

\textbf{Const} $\left[\;K(e_1,\ldots,e_h)\;\right]$\\[1em]
\begin{prooftree}
    \AxiomC{$e_1\downarrow n_1,\;\ldots,\;e_h\downarrow n_h$}
    \AxiomC{$\enc P\{^{n_1}/_{x_1},\ldots\}\cne \xrightarrow{\hat\alpha} Q$}
    \RightLabel{IH}
    \UnaryInfC{$P\{^{n_1}/_{x_1},\ldots\} \xrightarrow{\alpha} R \;\wedge\; \enc R\cne = Q$}
    \RightLabel{Const}
    \BinaryInfC{$K(e_1,\ldots,e_h) \xrightarrow{\alpha} R$}
\end{prooftree}
$\enc K(e_1,\ldots,e_h)\cne = K_{n_1,\ldots} \overset{\text{def}}{=} \enc P\{^{n_i}/_{x_i}\}\cne$.\\
The transition is on that encoding.\\
The IH recovers $P\{^{n_1}/_{x_1},\ldots\} \xrightarrow{\alpha} R$ with $\enc R\cne = Q$.\\
Taking $P' = R$ witnesses $K(e_1,\ldots,e_h) \xrightarrow{\alpha} P'$ and $\enc P'\cne = Q$.

\newpage

% ================================================================
% IMPLEMENTATION REMARK
% ================================================================

The interpreter developed for this project encodes input actions over a bounded interval
$[n_{\min}, n_{\max}] \subset \mathbb{N}$ rather than an infinite sum over all of $\mathbb{N}$.
Concretely, $\enc a(x).P\cne = \sum_{n=n_{\min}}^{n_{\max}} a_n.\enc P\{^n/_x\}\cne$
instead of $\sum_{n \in \mathbb{N}} a_n.\enc P\{^n/_x\}\cne$.
The proof above remains structurally identical: the \textbf{Input} base case still
applies the Sum rule and the $=$ step, and every inductive case goes through unchanged.
The only difference is that operational correspondence holds under the assumption that
all communicated values fall within the chosen interval; values outside it have no
corresponding branch in the encoding and therefore no matching CCS transition.

\end{document}
